\subsubsection{Représentation rectiligne des phases et du retour}
\label{mono-recta:desarrollo}

La représentation sur le segment \([0,10]\) conserve l'origine positionnelle d'APP et ordonne ses neuf marques intérieures. Dans ce diagramme, les entiers \(1,\ldots,9\) sont des représentants de phase. La position horizontale indique l'ordre des coordinatisations ; les étiquettes supérieures spécifient des opérations régionales. La figure \ref{mono-recta:figura} réunit ainsi la droite de départ, la fixation de l'axe d'auto-échelle, la distribution spéculaire des deux canaux et le retour avec mémoire (section~\ref{mono-ampl:conexion}).

% Recuperación de fig_nonadic_route_mini.tex de la síntesis académica.
% Los rótulos designan operaciones regionales; la suma se efectúa en F_3^6.
\begin{figure}[H]
\centering
\begin{tikzpicture}[x=1.12cm,y=1cm,font=\small,>=Latex,
  guide/.style={line width=.45pt,black!55}]
 \draw[line width=.8pt] (0,0)--(10,0);
 \foreach \k in {0,...,10}{
  \draw[guide] (\k,-.09)--(\k,.09);
  \node[below=2.5mm] at (\k,0) {$\k$};
 }
 \foreach \k in {1,...,9}{\fill (\k,0) circle (1.7pt);}
 \node[anchor=west,font=\footnotesize] at (0,3.05)
   {Segment générateur $[0,10]$ et neuf représentants de phase};
 \draw[guide] (5,.12)--(5,1.7);
 \node[above,align=center,font=\footnotesize] at (5,1.7)
   {fixation de l'axe\\$u^\varphi$};
 \draw[decorate,decoration={brace,amplitude=4pt},guide]
   (6,.65)--(8,.65);
 \node[above=5pt,align=center,font=\footnotesize] at (7,.65)
   {agrégat régional\\$u^e+u^\pi$};
 \draw[guide] (9,.12)--(9,1.7);
 \node[above,align=center,font=\footnotesize] at (9,1.7)
   {involution spéculaire\\$\pi\leftrightarrow e$};
 \draw[->,line width=.7pt] (9,-.7)
   .. controls (9,-1.55) and (1,-1.55) .. (1,-.7);
 \node[fill=white,inner sep=3pt,font=\footnotesize] at (5,-1.23)
   {retour $9\to1$;\quad $m\mapsto m+1$};
 \node[font=\footnotesize] at (5,-1.95)
   {$w_6\to w_{12}\to w_{18}\to w_{24}\to w_{30}\to R_{36}\to G_9$};
\end{tikzpicture}
\caption{Représentation rectiligne de la connexion nonadique. La marque
\(5\) situe la première saturation forte et l'axe régional d'autoéchelle ;
l'accolade \(6\)--\(8\) identifie l'agrégat ternaire
fixé dans chaque fibre ; la phase \(9\) sélectionne l'orientation de
clôture. Le retour conserve la phase et incrémente la mémoire. Les
abscisses numérotent les positions opératives, tandis que les étiquettes
\(u^\pi,u^e,u^\varphi\) désignent des blocs de \(\mathbb F_3^6\).}
\label{mono-recta:figura}
\end{figure}


La fixation de l'axe se formule dans la fibre d'une frontière donnée. Si
\(q=u^\pi+u^e+u^\varphi\) et
\(a=u^\pi+u^e-u^\varphi\), alors
\[
 u^\varphi=2(q-a),\qquad s=u^\pi+u^e=q-u^\varphi
 \quad\text{en }\mathbb F_3^6.
\]
Les données \((q,a)\) déterminent l'axe et l'agrégat ; la répartition ordonnée de \(s\) conserve la liberté spéculaire résolue par le prolongement. La cinquième phase possède une solution locale et une extension, avec la donnée de frontière \(c=111111\) et la signature résiduelle \(\rho_W=000\). C'est le sens de sa première saturation forte. Le symbole \(\varphi\) situé sur la marque \(5\) identifie l'axe régional qui intervient dans cette opération.

Les phases intermédiaires conservent cette décomposition au sein de chaque fibre,
tandis que leurs données de frontière évoluent. Le recensement complet des
multiplicités locales et des extensions de la branche distinguée est
\[
\begin{array}{c|rrrrrrrrr}
\text{phase}&1&2&3&4&5&6&7&8&9\\\hline
N_{\rm loc}&3&2&5&4&1&1&3&3&2\\
N_{\rm ext}&2&5&4&1&1&3&3&2&1
\end{array}
\]
La phase \(6\) combine unicité locale et trois prolongations ; la phase
\(7\) résout une fibre spéculaire triple ; la phase \(8\) synchronise
les trois recensements régionaux. La phase \(9\) conserve deux distributions
locales ayant le même axe et le même agrégat, dont l'une se
prolonge. Ces dénombrements expriment des fonctions différentes au sein d'une
même connexion et justifient les marques de la droite.

En particulier, pour les phases \(6,7,8\), les sommes régionales sont
respectivement \(012100\), \(101001\) et \(112000\). La conservation
représentée par l'accolade est une invariance par rapport à la répartition
spéculaire dans chaque fibre. L'évaluation réelle ultérieure applique les
lecteurs des cylindres compatibles aux histoires complètes ; la
somme interne représentée dans la figure appartient au domaine ternaire.

Enfin, l'arc \(9\to1\) représente le passage d'un tour au suivant. Les formules de la section \ref{mono-ampl:conexion} conservent le reste de phase et incrémentent le compteur de tours. Le diagramme rectiligne et le diagramme cyclique de la figure \ref{mono-ampl:fig-regional} expriment donc deux aspects complémentaires de la même holonomie : l'ordre des positions et la composition du retour sur l'état avec mémoire.
